% GENERATED FILE. Edit canonical lessons, then run npm run generate:books.
% canonical-source-hash: fnv1a64:6665f2e0fcdd0b14
% canonical-lessons: GU-C458-virah, GU-C458-ekat, GU-C458-sunu, GU-C458-tarasvu, GU-C458-sabharvu

\chapter{Missing Someone}
\label{ch:gu-missing-someone}
\hlchaptermodality{458}

\begin{chapteropening}
\textbf{By the end of this chapter:} \emph{I can say how I feel when someone is far away.}

Complete the last lesson of chapter 458: I can say how I feel when someone is far away.
\end{chapteropening}

\section[\texorpdfstring{\gu{વિરહ} (virah)}{વિરહ (virah)}]{\gu{વિરહ} (virah) — the pain of separation}

\label{lesson:GU-C458-virah}



\begin{quote}
Before the new one: say the Gujarati for envy, then the Gujarati for dislike.
\end{quote}

\subsection*{You'll want to know: \gu{વિરહ}}
\textbf{\gu{વિરહ}} — \emph{virah} — \textquotedblleft{}the pain of separation\textquotedblright{}.

\begin{quote}
\gu{પતિ} \gu{દૂર} \gu{ગયો} \gu{છે}, \gu{તેથી} \gu{પત્નીને} \gu{વિરહનું} \gu{દુઃખ} \gu{છે}.
\end{quote}

\begin{quote}
\emph{pati dūr gayo che, tethī patnīne virahnũ duḥkh che.}
\end{quote}

\begin{quote}
Her husband has gone far away, so the wife suffers the pain of separation.
\end{quote}

A masculine noun.

\begin{cousinweb}
\textbf{\gu{વિરહ}} (\emph{virah}) is a tatsama, Sanskrit \emph{viraha}, \textquotedblleft{}abandonment, parting, separation, especially of lovers\textquotedblright{}, as Monier-Williams records it.
\end{cousinweb}

\begin{grammarlens}[title={What you've built}]
One more word for describing what you lived and felt.
\end{grammarlens}

\subsection*{Your turn}
\begin{itemize}
\raggedright
  \item \emph{Say it:} \emph{virah}
  \item \emph{Say it:} \emph{virah}, once more
  \item \emph{Say it:} \emph{pati dūr gayo che, tethī patnīne virahnũ duḥkh che.} — Her husband has gone far away, so the wife suffers the pain of separation.
  \item \emph{Recall:} say the Gujarati for envy, then the Gujarati for dislike, then say \emph{virah} again
\end{itemize}

\begin{tcolorbox}[breakable,colback=teal!4,colframe=teal!35!black,title={Before you move on}]
What does \gu{વિરહ} mean? (\textquotedblleft{}The pain of separation\textquotedblright{}.) Say it once more.
\end{tcolorbox}
\section[\texorpdfstring{\gu{એકાંત} (ekā̃t)}{એકાંત (ekā̃t)}]{\gu{એકાંત} (ekā̃t) — solitude}

\label{lesson:GU-C458-ekat}



\begin{quote}
Before the new one: say the Gujarati for dislike, then the Gujarati for the pain of separation.
\end{quote}

\subsection*{You'll want to know: \gu{એકાંત}}
\textbf{\gu{એકાંત}} — \emph{ekā̃t} — \textquotedblleft{}solitude\textquotedblright{}.

\begin{quote}
\gu{દાદાને} \gu{સાંજે} \gu{એકાંત} \gu{ગમે} \gu{છે}.
\end{quote}

\begin{quote}
\emph{dādāne sā̃je ekā̃t game che.}
\end{quote}

\begin{quote}
Grandfather likes solitude in the evenings.
\end{quote}

A neuter noun.

\begin{cousinweb}
\textbf{\gu{એકાંત}} (\emph{ekā̃t}) is a tatsama, Sanskrit \emph{ekānta}, \textquotedblleft{}a lonely or retired place\textquotedblright{}, as Monier-Williams records it.
\end{cousinweb}

\begin{grammarlens}[title={What you've built}]
One more word for describing what you lived and felt.
\end{grammarlens}

\subsection*{Your turn}
\begin{itemize}
\raggedright
  \item \emph{Say it:} \emph{ekā̃t}
  \item \emph{Say it:} \emph{ekā̃t}, once more
  \item \emph{Say it:} \emph{dādāne sā̃je ekā̃t game che.} — Grandfather likes solitude in the evenings.
  \item \emph{Recall:} say the Gujarati for dislike, then the Gujarati for the pain of separation, then say \emph{ekā̃t} again
\end{itemize}

\begin{tcolorbox}[breakable,colback=teal!4,colframe=teal!35!black,title={Before you move on}]
What does \gu{એકાંત} mean? (\textquotedblleft{}Solitude\textquotedblright{}.) Say it once more.
\end{tcolorbox}
\section[\texorpdfstring{\gu{સૂનું} (sūnũ)}{સૂનું (sūnũ)}]{\gu{સૂનું} (sūnũ) — deserted; empty and silent}

\label{lesson:GU-C458-sunu}



\begin{quote}
Before the new one: say the Gujarati for the pain of separation, then the Gujarati for solitude.
\end{quote}

\subsection*{You'll want to know: \gu{સૂનું}}
\textbf{\gu{સૂનું}} — \emph{sūnũ} — \textquotedblleft{}deserted; empty and silent\textquotedblright{}.

\begin{quote}
\gu{દીકરીના} \gu{લગ્ન} \gu{પછી} \gu{ઘર} \gu{સૂનું} \gu{લાગે} \gu{છે}.
\end{quote}

\begin{quote}
\emph{dīkrīnā lagna pachī ghar sūnũ lāge che.}
\end{quote}

\begin{quote}
After the daughter's wedding the house feels empty.
\end{quote}

\begin{cousinweb}
\textbf{\gu{સૂનું}} (\emph{sūnũ}) is inherited from Sanskrit \emph{śūnya}, \textquotedblleft{}empty\textquotedblright{}; CDIAL lists the Gujarati word (CDIAL 12567).
\end{cousinweb}

\begin{grammarlens}[title={What you've built}]
One more word for describing what you lived and felt.
\end{grammarlens}

\subsection*{Your turn}
\begin{itemize}
\raggedright
  \item \emph{Say it:} \emph{sūnũ}
  \item \emph{Say it:} \emph{sūnũ}, once more
  \item \emph{Say it:} \emph{dīkrīnā lagna pachī ghar sūnũ lāge che.} — After the daughter's wedding the house feels empty.
  \item \emph{Recall:} say the Gujarati for the pain of separation, then the Gujarati for solitude, then say \emph{sūnũ} again
\end{itemize}

\begin{tcolorbox}[breakable,colback=teal!4,colframe=teal!35!black,title={Before you move on}]
What does \gu{સૂનું} mean? (\textquotedblleft{}Deserted; empty and silent\textquotedblright{}.) Say it once more.
\end{tcolorbox}
\section[\texorpdfstring{\gu{તરસવું} (tarasvũ)}{તરસવું (tarasvũ)}]{\gu{તરસવું} (tarasvũ) — to long for; to crave}

\label{lesson:GU-C458-tarasvu}



\begin{quote}
Before the new one: say the Gujarati for solitude, then the Gujarati for deserted.
\end{quote}

\subsection*{You'll want to know: \gu{તરસવું}}
\textbf{\gu{તરસવું}} — \emph{tarasvũ} — \textquotedblleft{}to long for; to crave\textquotedblright{}.

\begin{quote}
\gu{દાદી} \gu{દીકરાને} \gu{જોવા} \gu{તરસે} \gu{છે}.
\end{quote}

\begin{quote}
\emph{dādī dīkrāne jovā tarse che.}
\end{quote}

\begin{quote}
Grandmother is longing to see her son.
\end{quote}

Past (he …): \gu{તે} \gu{ઘરે} \gu{જવા} \gu{તરસતો} \gu{હતો}, \emph{te ghare javā tarasto hato}.

\begin{cousinweb}
\textbf{\gu{તરસવું}} (\emph{tarasvũ}) goes back to a reconstructed \emph{*tṛṣati}, \textquotedblleft{}is thirsty\textquotedblright{}; CDIAL lists the Gujarati verb as \textquotedblleft{}is thirsty\textquotedblright{} (CDIAL 5935), and Gujarati also uses it for longing, a thirst of the heart.
\end{cousinweb}

\begin{grammarlens}[title={What you've built}]
One more word for describing what you lived and felt.
\end{grammarlens}

\subsection*{Your turn}
\begin{itemize}
\raggedright
  \item \emph{Say it:} \emph{tarasvũ}
  \item \emph{Say it:} \emph{tarasvũ}, once more
  \item \emph{Say it:} \emph{dādī dīkrāne jovā tarse che.} — Grandmother is longing to see her son.
  \item \emph{Recall:} say the Gujarati for solitude, then the Gujarati for deserted, then say \emph{tarasvũ} again
\end{itemize}

\begin{tcolorbox}[breakable,colback=teal!4,colframe=teal!35!black,title={Before you move on}]
What does \gu{તરસવું} mean? (\textquotedblleft{}To long for; to crave\textquotedblright{}.) Say it once more.
\end{tcolorbox}
\section[\texorpdfstring{\gu{સાંભરવું} (sā̃bharvũ)}{સાંભરવું (sā̃bharvũ)}]{\gu{સાંભરવું} (sā̃bharvũ) — to come to mind; to be missed}

\label{lesson:GU-C458-sabharvu}



\begin{quote}
Before the new one: say the Gujarati for deserted, then the Gujarati for to crave.
\end{quote}

\subsection*{You'll want to know: \gu{સાંભરવું}}
\textbf{\gu{સાંભરવું}} — \emph{sā̃bharvũ} — \textquotedblleft{}to come to mind; to be missed\textquotedblright{}.

\begin{quote}
\gu{રજામાં} \gu{મને} \gu{મારું} \gu{ગામ} \gu{ઘણું} \gu{સાંભરે} \gu{છે}.
\end{quote}

\begin{quote}
\emph{rajāmā̃ mane mārũ gām ghaṇũ sā̃bhre che.}
\end{quote}

\begin{quote}
During the holidays I miss my village a lot.
\end{quote}

Past (it …): \gu{તે} \gu{ગામ} \gu{ઘણું} \gu{સાંભર્યું}, \emph{te gām ghaṇũ sā̃bharyũ}.

\begin{cousinweb}
\textbf{\gu{સાંભરવું}} (\emph{sā̃bharvũ}) is inherited from Sanskrit \emph{saṃsmarati}, \textquotedblleft{}remembers\textquotedblright{}; CDIAL lists the Gujarati verb, \textquotedblleft{}to remember\textquotedblright{} (CDIAL 13057).
\end{cousinweb}

\begin{grammarlens}[title={What you've built}]
One more word for describing what you lived and felt.
\end{grammarlens}

\subsection*{Your turn}
\begin{itemize}
\raggedright
  \item \emph{Say it:} \emph{sā̃bharvũ}
  \item \emph{Say it:} \emph{sā̃bharvũ}, once more
  \item \emph{Say it:} \emph{rajāmā̃ mane mārũ gām ghaṇũ sā̃bhre che.} — During the holidays I miss my village a lot.
  \item \emph{Recall:} say the Gujarati for deserted, then the Gujarati for to crave, then say \emph{sā̃bharvũ} again
\end{itemize}

\begin{tcolorbox}[breakable,colback=teal!4,colframe=teal!35!black,title={Before you move on}]
What does \gu{સાંભરવું} mean? (\textquotedblleft{}To come to mind; to be missed\textquotedblright{}.) Say it once more.
\end{tcolorbox}
